\chapter{Research Design \& Methods}%
\label{ch:methods}
This chapter introduces the research approach and general methodology that was used in this study.

\section{Research approach}
This study follows an evaluative and design-oriented software-engineering approach. It first investigates the activities and challenges involved in experiments with anomaly-detection (AD) and root-cause-analysis (RCA) methods. It then derives workflow requirements, translates these requirements into a software design, and evaluates a working implementation.

CloudSentinel is used as the main case and implementation vehicle. Studying its original workflow provides concrete evidence about the work needed to move from Kubernetes telemetry to training, anomaly detection, and RCA. Redesigning the framework then makes it possible to investigate how the identified activities can be supported as one connected experimental lifecycle. The research questions are therefore broader than CloudSentinel, while the claims about the implemented solution remain bounded by this case.

\section{Research steps}
The study consists of four related phases.

\subsection{Phase 1: activities and practical challenges}

[TODO]

\subsection{Phase 2: workflow requirements}

[TODO]

\subsection{Phase 3: workflow design and implementation}

[TODO]

\subsection{Phase 4: evaluation of the implemented workflow}

[TODO]

\section{Data \& validity}
[TODO]










% Gebruik van synthetische data i.p.v. echte klantdata (privacy/beleid).
% Beperkingen van synthetische workloads, etc.

% ---------------------------------------
% In this chapter, you described the design of your research: which methods and procedures you used as a step towards answering the research questions.
% A well-written methods chapter highlights the plausibility of your research methods.
% You should clarify to the reader \emph{why} you have chosen specific methods and how they are justified for your research.
% To garner more credibility, you can include the pitfalls and difficulties associated with your approach.

% Depending on the nature of your research, you may, for example, want to describe the exact procedure you use to collect the data, how it is analysed, and how you plan to validate the data and results. 
% On the other hand, when you develop a new framework, you want to describe in detail the requirements and the design, and the motivation behind them (\ie{} the decision you made and why).
% Make sure to discuss all ethical aspects of your research.
% This is more important when dealing with (datasets of) people or (trained) AI models.
